% MT034 / explanation — canonical editable content.
题目的诱惑是把“样本无限多”理解成“所有问题都能学会”。更多数据能让我们更准确地知道各类怎样分布，却不能让本来重叠的两类自动分开。


补充例子：设类0与类1的特征分布完全相同，而且各占一半。那么看到任何\(x\)，两个类别的后验都是0.5。你选0，会错掉标签为1的那一半；选1，也会错一半。无限多数据只是让你更确定这个50\%与50\%的事实，最小错误率仍是50\%。


一般地，选后验最大的类别后，在给定\(x\)处的错误概率是
\[
P(\text{error}\mid x)=1-\max_c P(y=c\mid x).
\]


例如最高后验是0.8，即使做了当前最合理的决定，仍有0.2概率遇到另一类。把这些无法区分造成的错误在所有输入上平均，得到Bayes错误率。只要重叠区域有正概率，它就可能大于0。


题目说的是训练误差，但分类器被限定为按正确分布建立的Bayes分类器，不能改成逐条记住标签的查表程序。要保证零错误，还需要几乎每个可能遇到的输入都只对应一个类别，即后验概率只能是0或1；原题没有给出这个条件，故一般为False。


\textbf{出处：} 2021B*期中 Q8，原卷第3页（10分）


\textbf{答案依据：} 2021B配套解答 / 2021B配套答案副本 Q8随卷解答。
